\setcounter{ex}{0}
\setcounter{paragraph}{0}
\PhanI
  \begin{ex}
\immini[thm]{Cho đồ thị li độ - thời gian của một dao động điều hòa như hình vẽ. Quãng đường vật đi được từ thời điểm $0,5\ s$ đến $1,5\ s$ là
\choice[2]
 	{\True $2\ cm$}
 	{$1\ cm$}
 	{$4\ cm$}
 	{$2,5\ cm$}}{
\begin{hinhanh}{7}
\begin{tikzpicture}[very thick,>=stealth']
\pgfplotsset{width=7cm,height=4.5cm,compat=1.4}
\begin{axis}[
axis lines=middle, xmin=0, xmax=2*pi, xstep=1,
ymin=-4.5, ymax=5, ystep=0.1,
grid=major,%Vẽ chế độ lưới theo trục tọa độ Oxy,
x grid style={dashed, black},
axis line style={very thick,Mapcolor}, % <--- dòng thêm vào
xlabel=$t~(s)$,
xlabel style={above},
xtick ={1,2,3,4,5,6,7,8},
xticklabels={},
y grid style={dashed, black},
extra y ticks={0},
extra y tick label={0},
ylabel=$x~(cm)$,
ylabel style={above},
ytick={-4,0, 4},
yticklabels={$-2$,,2},]
\addplot[line width=1.2pt, domain=0:1.7*pi, samples=400, Mapcolor]{4*cos(deg(0.5*pi*x-pi/2))};
\addplot[Mapcolor,mark=*,mark options={fill=white,mark size=1.5pt}, nodes near coords={0,5}]
coordinates {(1, 0)};
\addplot[Mapcolor,mark=*,mark options={fill=white,mark size=1.5pt}, nodes near coords={1,5}]
coordinates {(3, 0)};
\addplot[Mapcolor,mark=*,mark options={fill=white,mark size=1.5pt}, nodes near coords={2,5}]
coordinates {(5, 0)};
\end{axis}
\end{tikzpicture}
\end{hinhanh}}
 	\loigiai{
 		Từ đồ thị, ta thấy biên độ dao động của vật là $A = 1\ cm$. Chu kì dao động là $T = 2 s$.
 		Tại thời điểm $t_1 = 0,5 s$, li độ của vật là $x(0,5) = 1\ cm$ (vật ở vị trí biên dương).
 		Tại thời điểm $t_2 = 1,5 s$, li độ của vật là $x(1,5) = -1\ cm$ (vật ở vị trí biên âm).
 		Khoảng thời gian $\Delta t = t_2 - t_1 = 1,5 s - 0,5 s = 1 s = \dfrac{T}{2}$.
 		Trong nửa chu kì, vật đi được quãng đường bằng $2A = 2 . 1 cm = 2\ cm$.
 	}
 \end{ex}
  %\loai{Tính quãng đường vật đi được trong quãng thời gian $\Delta t=nT$}
  \begin{ex}%[3P1KD-1][Loai1] 
  	Một chất điểm dao động điều hoà với chu kì $T=\dfrac{\pi}{10}\ s$ và biên độ $A = 6$ cm. Chọn gốc thời gian $t = 0$ lúc vật qua vị trí cân bằng. Quãng đường vật đi được trong $10\pi$ s đầu tiên là
  	\choice
  	{9 m}
  	{\True 24 m}
  	{6 m}
  	{1 m}
  	\loigiai{
  		\begin{itemize}
  			\item Ta có:
  			$t = 100T \Rightarrow S = 100.4A = 400A = 24\ m$ (1T đi được 4A).
  			
  		\end{itemize}
  		
  	}
  	%\haidong{8}
  \end{ex} \haidong{20}
  %\loai{Tính quãng đường vật đi được trong quãng thời gian $\Delta t=(2n+1)\dfrac{T}{2}$}
  \begin{ex}%[3P1BD-1][Loai1] 
  	Một chất điểm dao động điều hoà với tần số góc $\omega=20\ rad/s$ và biên độ $A = 6$ cm. Chọn gốc thời gian lúc vật đi qua vị trí cân bằng. Quãng đường vật đi được trong $0,05\pi$ s đầu tiên là
  	\choice
  	{24 cm}
  	{9 cm}
  	{6 cm}
  	{\True 12 cm}
  	\loigiai{
  		\begin{itemize}
  			\item Trong $0,05\pi = \dfrac{T}{2}$ Vật đi được quãng đường $2A = 12$ cm dù cho gốc thời gian là ở lúc nào.
  			\item 
  		\end{itemize}
  	}
  	%\haidong{8}
  \end{ex} \haidong{20}
  \loai{Quãng đường vật đi được sau thời gian $\Delta t$ kể từ thời điểm ban đầu - không cần tách thời gian}
  \begin{ex}%[3P1KD-1][Loai1]
  	Một vật nhỏ dao động điều hòa có biên độ A, chu kì dao động T , ở thời điểm ban đầu  $t_0=0$  vật đang ở vị trí biên. Quãng đường mà vật đi được từ thời điểm ban đầu đến thời điểm $t=\dfrac{T}{3}$ là
  	\choice
  	{\True $\dfrac{3A}{2}$}
  	{$\dfrac{A}{3}$}
  	{$\dfrac{2A}{3}$}
  	{A}
  	\loigiai{
  		\begin{itemize}
  			\item Sau $\dfrac{T}{3}$, nếu vật ở biên dương sẽ về $-\dfrac{A}{2}$, nếu vật ở biên âm sẽ tới $\dfrac{A}{2}$ $\Rightarrow$ Quãng đường đi được là $\dfrac{3A}{2}$.
  			\item 
  			\item 
  		\end{itemize}
  	}
  	%\haidong{7}
  \end{ex} \haidong{20}
  \loai{Quãng đường vật đi được sau thời gian $\Delta t$ kể từ thời điểm ban đầu - cho phương trình dao động}
  \begin{ex}%[3P1KD-1][Loai1] 
  	Một vật dao động điều hoà với phương trình $x = 4\cos(\pi t+ \dfrac{\pi}{4})$ cm. Sau 4,5 s kể từ thời điểm đầu tiên vật đi được đoạn đường là
  	\choice
  	{34 cm}
  	{36 cm}
  	{\True $32 + 4\sqrt{2}$ cm}
  	{$32 + 2\sqrt{2}$ cm}
  	\loigiai{
  		\begin{itemize}
  			\item $t=0$ vật đang ở li độ $x = 2\sqrt{2}$ theo chiều âm.
  			\item $\Delta t = 4,5 = \dfrac{T}{4}+2T.$
  			\item $t=\dfrac{T}{4}$ vật ở li độ $x = -2\sqrt{2}$ theo chiều âm.
  			\item Quãng đường vật đi được trong $\dfrac{T}{4}$ s đầu tiên là $S_1 = 2.2\sqrt{2} = 4\sqrt{2}\ cm$.
  			\item Quãng đường đi được là $S = S_1 + 8A = 32 + 4\sqrt{2}\ cm$.
  		\end{itemize}
  	}
  	%\haidong{7}
  \end{ex} \haidong{20}
  \loai{Cho quãng đường viết phương trình dao động}
  \begin{ex}%[3P1KD-1][Loai1]
  	Một chất điểm dao động điều hòa với tần số góc $\omega=20\pi\ rad/s$. Biết rằng trong quãng thời gian 0,5 s vật nhỏ đi được tổng quãng đường bằng 50 cm. Và tại thời điểm $t = 0$, vật ở biên dương. Lấy gần đúng $\pi^2 = 10$. Phương trình dao động của vật là
  	\choice 
  	{$x = 2,5\cos(20\pi t + \pi)$ cm}
  	{$x = 5\cos(20\pi t)$ cm}
  	{$x = 5\cos(20\pi t + \pi)$ cm}
  	{\True $x = 2,5\cos(20\pi t)$ cm}
  	\loigiai{
  		\begin{itemize}
  			\item Ta có:
  			$T = \dfrac{2\pi}{20\pi} = 0,1\ s
  			\Rightarrow t = 0,5\ s = 5T \Rightarrow s = 5.4A = 50 \Rightarrow A = 2,5\ cm$. 
  			\item Phương trình dao động của vật là: $x = 2,5\cos(20\pi t + \varphi)\ cm$.
  			\item Tại $t = 0$, $x = 2,5 \Rightarrow 2,5\cos\varphi = 2,5 \Rightarrow \cos\varphi = 1 \Rightarrow \varphi = 0$.
  			\item Phương trình dao động của vật là $x = 2,5\cos(20\pi t)$ cm.
  		\end{itemize}
  	}
  	%\haidong{7}
  \end{ex} \haidong{20}
  
  \loai{Tìm thời gian vật đi được quãng đường $s < 2A$}
  \begin{ex}%[3P1KD-2][Loai1] 
  	Một vật dao động điều hoà theo phương trình $x = 4\cos(8\pi t - \dfrac{2\pi}{3})$ cm. Thời gian vật đi được quãng đường $s = (2+2\sqrt{2})$ cm kể từ lúc vật bắt đầu dao động là
  	\choice
  	{\True $\dfrac{5}{96}$ s}
  	{$\dfrac{1}{96}$ s}
  	{$\dfrac{29}{96}$ s}
  	{$\dfrac{25}{96}$ s}
  	\loigiai{
  		\begin{itemize}
  			\item Ta có ban đầu vật đang ở li độ $x=-2$ cm theo chiều dương.
  			\item Sau khi đi được quãng đường $S = 2 + 2\sqrt{2}\ cm$ vật đang ở li độ: 
  			\begin{align*}
  				x = \dfrac{A\sqrt{2}}{2} \Rightarrow t = \dfrac{T}{12}+ \dfrac{T}{8} = \dfrac{5T}{24} = \dfrac{5}{96}\ s.
  			\end{align*}
  			
  		\end{itemize}
  	}
  	%\haidong{7}
  \end{ex} \haidong{20}
  \loai{Tìm thời gian vật đi được quãng đường $2A<s < 4A$}
  \begin{ex}%[3P1KD-2][Loai1]
  	Cho chất điểm đang dao động điều hòa trên trục Ox quanh vị trí cân bằng là gốc tọa độ O với phương trình li độ $x = 4\cos(2\pi t - \dfrac{\pi}{3})$ cm. Tính từ thời điểm ban đầu, thời gian để chất điểm đi được quãng đường bằng 14 cm là
  	\choice
  	{\True $\dfrac{11}{12}$ s}
  	{$\dfrac{5}{6}$ s}
  	{$\dfrac{2}{3}$ s}
  	{$\dfrac{6}{7}$ s}
  	\loigiai{
  		\immini{
  			\begin{itemize}
  				\item 
  				Ta có: $s = 14\ cm = 2 + 12$.
  				\item Ban đầu chất điểm ở vị trí có li độ 2 cm ($M_0$), để đi được 14 cm thì nó phải quay đến điểm $M_1$
  				$\Delta \varphi = \dfrac{11\pi}{6} \Rightarrow t = \dfrac{11T}{12} = \dfrac{11}{12}\ s$.
  				
  			\end{itemize}
  		}{\begin{tikzpicture}[scale=1,thick,>=stealth']
  				\tkzInit[ymin=-3,ymax=3,xmin=-3,xmax=3]\tkzClip
  				\tkzDefPoints{-2.5/0/m, 2.5/0/n, 0/0/o, 0/2.5/p, 0/-2.5/q}
  				\tkzDefShiftPoint[o](-90:2){m'}
  				\tkzDefShiftPoint[o](-60:2){0}
  				\tkzDefPointBy[projection=onto m--n](0) \tkzGetPoint{h}
  				\draw(o)circle(2cm);
  				\tkzDrawSegments[dashed](h,0)
  				\tkzDrawSegments(p,q)
  				\tkzDrawSegments[->](m,n)
  				\tkzDrawSegments[blue,->](o,0 o,m')
  				\tkzDrawPoints[size=3](o,h,0,m')
  				\tkzMarkAngles[size=0.7cm,arrows=->](0,o,m')
  				\node at (0)[below right]{$M_0$};
  				\node at (m')[below left]{$M_1$};
  				\node at (2,0)[below right]{$4$};
  				\node at (h)[below right]{$2$};
  		\end{tikzpicture}}
  	}
  	%\haidong{7}
  \end{ex} \haidong{20}
  \begin{ex}
  	Một chất điểm dao động điều hòa với phương trình $x = 4 \cos (4 \pi t)\ cm$. Quãng đường vật đi được trong thời gian $2,875\ s$ kể từ lúc $t = 0$ là
  	\choice
  	{$88\ cm$}
  	{\True $92\ cm$}
  	{$96\ cm$}
  	{$100\ cm$}
  	\loigiai{
  		Ta có $T = \dfrac{2\pi}{\omega} = \dfrac{2\pi}{4\pi} = 0,5\ s$\\
  		\(\Rightarrow\) Số chu kì đã dao động: \(n = \dfrac{2,875}{0,5} = 5,75\)\\
  		\(\Rightarrow\) Quãng đường: \(S = 5. 4A + 3A = 20A + 3A = 23A = 23. 4 = 92\ cm\)
  	}
  \end{ex}
  
  %\loai{Tổng hợp}
  \begin{ex}%[3P1GD-1][Loai1]
  	Cho chất điểm dao động điều hòa với biên độ bằng 8 cm. Thời điểm ban đầu, $t=0$, chất điểm đi qua vị trí li độ $-4$ cm theo chiều âm. Trong 0,5 s đầu tiên chất điểm đi được quãng đường dài ($44-4\sqrt{3}$) cm. Tốc độ cực đại của chất điểm trong quá trình dao động là
  	\choice
  	{$8\pi$ cm/s}
  	{$36\pi$ cm/s}
  	{\True $40\pi$ cm/s}
  	{$20\pi$ cm/s}
  	\loigiai{
  		\immini{\begin{itemize}
  				\item 
  				Ta biểu diễn vị trí của chất điểm tại thời điểm ban đầu bằng điểm $M_0$ trên đường tròn lượng giác như hình vẽ.
  				\item Ta có: $s = 44-4\sqrt{3} = 32+4+8-4\sqrt{3}$.
  				\item Trong 0,5 s, chất điểm từ vị trí $M_0$ quay đúng 1 vòng, trở lại $M_0$ sau đó quay thêm góc $\dfrac{\pi}{2}$ và đến $M_1\Rightarrow t = T+\dfrac{T}{4} = 5\dfrac{T}{4} = 0,5 \Rightarrow T = 0,4 \Rightarrow \omega = 5\pi\ rad/s\Rightarrow v_{\max} = 5\pi.8 = 40 \pi\ cm/s$.
  		\end{itemize}}{\begin{tikzpicture}[scale=1,thick,>=stealth']
  				\tkzInit[ymin=-3,ymax=3,xmin=-3,xmax=3]\tkzClip
  				\tkzDefPoints{-2.5/0/m, 2.5/0/n, 0/0/o, 0/2.5/p, 0/-2.5/q}
  				\tkzDefShiftPoint[o](120:2){0}
  				\tkzDefShiftPoint[o](210:2){m'}
  				\tkzDefPointBy[projection=onto m--n](0) \tkzGetPoint{h}
  				\tkzDefPointBy[projection=onto m--n](m') \tkzGetPoint{h'}
  				\draw(o)circle(2cm);
  				\tkzDrawSegments[dashed](m',h' 0,h)
  				\tkzDrawSegments(p,q)
  				\tkzDrawSegments[->](m,n)
  				\tkzDrawSegments[blue,->](o,0 o,m')
  				\tkzDrawPoints[size=3](o,h,0,m',h')
  				\tkzMarkAngles[size=0.7cm,arrows=->](0,o,m')
  				\node at (0)[above left]{$M_0$};
  				\node at (m')[below left]{$M'$};
  				\node at (2,0)[below right]{$8$};
  				\node at (h)[below]{$-4$};
  				\node at (h')[above]{$-4\sqrt{3}$};
  		\end{tikzpicture}}
  	}
  	%\haidong{9}
  \end{ex} \haidong{20}
  

  
  
  %\loai{Tốc độ trung bình trong khoảng thời gian từ $t_1$ đến $t_2$}
  
  
  \begin{ex}%[3P1KD-3][Loai1]
  	Cho chất điểm M đang chuyển động tròn đều trên một quỹ đạo tròn có bán kính bằng 30 cm, với tốc độ góc là $\pi$ rad/s. Gọi P là hình chiếu của điểm M xuống một đường thẳng đi qua tâm và nằm trong mặt phẳng quỹ đạo. Tốc độ chuyển động trung bình của điểm P trong quãng thời gian 3 s là
  	\choice
  	{20 cm/s}
  	{\True 60 cm/s}
  	{$\dfrac{20}{3}$ cm/s}
  	{40 cm/s}
  	\loigiai{
  		\begin{itemize}
  			\item Trong khoảng thời gian 3 s chất điểm quay được góc: $\Delta \varphi = 3\pi = 2\pi + \pi\Rightarrow  \Delta s = 4A+2A= 180$ cm.
  			\item Tốc độ trung bình bằng: $180:3 = 60$ cm/s.
  		\end{itemize}
  	}
  	%\haidong{7}
  \end{ex} \haidong{20}
  \begin{ex}%[3P1KD-3][Loai1]
  	Cho một chất điểm dao động điều hòa theo phương trình $x = A\cos(\pi t + \dfrac{\pi}{4})$, x tính bằng cm và t tính bằng s. Tính từ thời điểm ban đầu, tốc độ trung bình của vật trong 3,5 s chuyển động là 6,36 cm/s. Biên độ dao động của chất điểm là
  	\choice
  	{6 cm}
  	{\True 3 cm}
  	{4 cm}
  	{5 cm}
  	\loigiai{
  		\immini{
  			\begin{itemize}
  				\item $T = \dfrac{2\pi}{\pi} = 2\ s$.
  				\item Ta có: $t = 3,5\ s = 1.2 + 1,5$.
  				\item Ta biểu diễn vị trí ban đâu của vật bằng điểm $M_0$ trên đường tròn lượng giác như hình vẽ.
  				\item Trong khoảng thời gian 2 s, vật quay được góc $2\pi$, tương ứng đi được quãng đường 4A và lại quay về vị trí $M_0$.
  				\item Trong khoảng thời gian 1,5 giây tiếp theo, vật quay thêm được 1 góc $\dfrac{3\pi}{2}$, đến vị trí $M_1$ trên đường tròn lượng giác.
  				\item Quãng đường chất điểm đi được trong 3,5 s là $s = 4A \dfrac{A}{\sqrt{2}}+2A+ \dfrac{A}{\sqrt{2}} = 6A+A\sqrt{2} = 2,5.6,36\Rightarrow  A = 3$ cm.
  			\end{itemize}
  		}{\begin{tikzpicture}[scale=1,thick,>=stealth']
  				\tkzInit[ymin=-3,ymax=3,xmin=-3,xmax=3]\tkzClip
  				\tkzDefPoints{-2.5/0/m, 2.5/0/n, 0/0/o, 0/2.5/p, 0/-2.5/q}
  				\tkzDefShiftPoint[o](45:2){0}
  				\tkzDefShiftPoint[o](-45:2){m'}
  				\tkzDefPointBy[projection=onto m--n](m') \tkzGetPoint{h}
  				\draw(o)circle(2cm);
  				\tkzDrawSegments[dashed](m',0)
  				\tkzDrawSegments(p,q)
  				\tkzDrawSegments[->](m,n)
  				\tkzDrawSegments[blue,->](o,0 o,m')
  				\tkzDrawPoints[size=3](o,h,0,m')
  				\tkzMarkAngles[size=0.7cm,arrows=->](0,o,m')
  				\node at (0)[above right]{$M_0$};
  				\node at (m')[below right]{$M_1$};
  				\node at (2,0)[below right]{$A$};
  				\node at (h)[below]{$\dfrac{A}{\sqrt{2}}$};
  		\end{tikzpicture}}
  	}
  	%\haidong{8}
  \end{ex} \haidong{20}
  %\loai{Tốc độ trung bình trong khoảng thời gian giữa hai lần liên tiếp vật qua li độ $x$}
  
  \begin{ex}%[3P1KD-3][Loai1] 
  	Một chất điểm dao động điều hòa với biên độ 5 cm và tần số 2 Hz. Tốc độ trung bình giữa hai lần liên tiếp chất điểm đi qua vị trí có li độ bằng 0 là
  	\choice
  	{20 cm/s}
  	{10 cm/s}
  	{\True 40 cm/s}
  	{30 cm/s}
  	\loigiai{
  		\begin{itemize}
  			\item Ta có: $T = \dfrac{1}{f} = 0,5\ s$.
  			\item Giữa hai lần liên tiếp chất điểm đi qua vị trí có li độ bằng 0 thì: 
  			\begin{align*}
  				t = \dfrac{T}{2} = 0,25\ s,\ S = 2A = 10\ cm \Rightarrow  v_{tb} = \dfrac{S}{t} = \dfrac{10}{0,25} = 40\ cm/s.
  			\end{align*}
  	\end{itemize}}
  	%\haidong{7}
  \end{ex} \haidong{20}
  
  
  %\loai{Tốc độ trung bình trong khoảng thời gian vật đi từ $x_1$ đến $x_2$}
  
  \begin{ex}%[3P1KD-3][Loai1] 
  	Cho dao động điều hòa với $f = 2\ Hz$ và $v_{\max} = 16\pi$ cm/s. Tìm tốc độ trung bình từ lúc chất điểm qua li độ 2 cm theo chiều âm đến lúc đi qua li độ $-2$ cm lần thứ hai.
  	\choice
  	{\True 32 cm/s}
  	{16 cm/s}
  	{24 cm/s}
  	{18 cm/s}
  	\loigiai{
  		\immini{
  			\begin{itemize}
  				\item 
  				Ta có: $T = \dfrac{1}{f} = 0,5$ s, $A = \dfrac{v_{\max}}{\omega} = \dfrac{16\pi}{2.2\pi} = 4\ cm$.
  				\item Vị trí có li độ $2\ cm$ và $-2\ cm$ được biểu diễn trên đường tròn lượng giác như hình vẽ. 
  				\item Ban đầu chất điểm qua li độ 2 cm theo chiều âm nên nó ở vị trí (3), khi nó đi qua vị trí li độ $-2$ cm lần thứ 2 thì nó ở vị trí $(2)$ nên quãng đường vật đi được là $2 + 4 + 2 = 8$ cm.
  				\item Thời gian để vật đi từ vị trí $(3)$ tới vị trí $(2)$ là $\dfrac{T}{2} = 0,25$ s
  				$\Rightarrow  v_{tb} = \dfrac{8}{0,25} = 32\ cm/s.$
  			\end{itemize}
  		}{\begin{tikzpicture}[scale=1,thick,>=stealth']
  				\tkzInit[ymin=-3,ymax=3,xmin=-3,xmax=3]\tkzClip
  				\tkzDefPoints{-2.5/0/m, 2.5/0/n, 0/0/o, 0/2.5/p, 0/-2.5/q}
  				\tkzDefShiftPoint[o](120:2){1}
  				\tkzDefShiftPoint[o](-120:2){2}
  				\tkzDefShiftPoint[o](60:2){3}
  				\tkzDefPointBy[projection=onto m--n](1) \tkzGetPoint{h}
  				\tkzDefPointBy[projection=onto m--n](3) \tkzGetPoint{h'}
  				\draw(o)circle(2cm);
  				\tkzDrawSegments[dashed](h',3 1,2)
  				\tkzDrawSegments(p,q)
  				\tkzDrawSegments[->](m,n)
  				\tkzDrawSegments[blue,->](o,1 o,2 o,3)
  				\tkzDrawPoints[size=3](o,1,h,2,3,h')
  				\node at (1)[above left]{$(1)$};
  				\node at (2)[below left]{$(2)$};
  				\node at (3)[above right]{$(3)$};
  				\node at (2,0)[below right]{$4$};
  				\node at (h')[below]{$3$};
  				\node at (h)[below left]{$-3$};
  		\end{tikzpicture}}
  	}
  	%\haidong{7}
  \end{ex} \haidong{20}
  %\loai{Tốc độ trung bình trong một chu kì}
  \begin{ex}%[3P1BD-3][Loai1]
  	Một vật dao động điều hòa có độ lớn vận tốc cực đại là 15,7 cm/s. Lấy $\pi =3,14$. Tốc độ trung bình của vật trong một chu kì dao động là
  	\choice
  	{20 cm/s}
  	{\True 10 cm/s}
  	{0 cm/s}
  	{15 cm/s}
  	\loigiai{
  		Tốc độ trung bình: $v_{tb(T)}=\dfrac{4A}{T}=\dfrac{2\omega A}{\pi}=\dfrac{2v_{\max }}{\pi} = 10$ cm/s.
  	}
  	%\haidong{5}
  \end{ex} \haidong{20}
   
  %\loai{Cho khoảng thời gian thỏa mãn tốc độ và tốc độ trung bình. Tìm các đại lượng}
  \begin{ex}%[3P1KD-3][Loai1] 
  	Cho một chất điểm đang dao động điều hòa với tần số bằng 1 Hz, biết trong mỗi chu kì dao động, khoảng thời gian mà tốc độ chuyển động của vật nhỏ hơn $5\pi\sqrt{2}$ cm/s là $\dfrac{1}{2}$ s. Tốc độ trung bình của chất điểm trong một chu kì là
  	\choice
  	{$20\pi$ cm/s}
  	{\True 20 cm/s}
  	{40 cm/s}
  	{$10\pi$ cm/s}
  	\loigiai{
  		\begin{itemize}
  			\item Ta có $\dfrac{1}{2}\ s = \dfrac{T}{2} \Rightarrow $ dùng đường tròn đơn vị cho vận tốc ta dễ dàng suy ra:
  			\begin{align*}
  				\dfrac{v_{\max}}{\sqrt{2}} = 5\pi \sqrt{2} \Rightarrow  v_{\max} = 10\pi\ cm/s \Rightarrow  A = \dfrac{v_{\max}}{\omega} = 5\ cm.
  			\end{align*}
  			\item Tốc độ trung bình của chất điểm trong 1 chu kì: $v_{tb}=\dfrac{4A}{T}=20$ cm/s.
  	\end{itemize}}
  	%\haidong{5}
  \end{ex} \haidong{20}
  %\loai{Tốc độ trung bình giữa hai thời điểm có li độ và gia tốc cho trước}
  \begin{ex}%[3P1KD-4][Loai1] 
  	Một vật nhỏ dao động điều hòa theo một quỹ đạo thẳng dài 20 cm với tần số 0,5 Hz. Từ thời điểm vật qua vị trí có li độ 5 cm theo chiều dương đến khi gia tốc của vật đạt giá trị cực tiểu lần thứ hai, vật có tốc độ trung bình là
  	\choice
  	{27,0 cm/s}
  	{26,7 cm/s}
  	{\True 19,29 cm/s}
  	{27,3 cm/s}
  	\loigiai{
  		\immini{\begin{itemize}
  				\item 
  				Ta có: $2A = 20 \Rightarrow  A = 10\ cm,\ T = \dfrac{1}{f} = 2\ s$.
  				\item Ta biểu diễn vị trí có li độ 5 cm theo chiều dương bằng điểm $M_0$ trên hình vẽ.
  				\item Gia tốc của vật có giá trị cực tiểu tại biên dương (vị trí điểm M' trên đường tròn).
  				\item Từ hình vẽ, ta suy ra quãng đường vật đi từ thời điểm vật qua vị trí có li độ 5 cm theo chiều dương đến khi gia tốc của vật đạt giá trị cực tiểu lần thứ hai là $s = 5 + 10.4 = 45$ cm.
  				\item Thời gian để vật qua vị trí có li độ 5 cm theo chiều dương đến khi gia tốc của vật đạt giá trị cực tiểu lần thứ hai là $\dfrac{T}{6} + T = \dfrac{7T}{6} = \dfrac{7}{3}\ s$.
  		\end{itemize}}{\begin{tikzpicture}[scale=1,thick,>=stealth']
  				\tkzInit[ymin=-3,ymax=3,xmin=-3,xmax=3]\tkzClip
  				\tkzDefPoints{-2.5/0/m, 2.5/0/n, 0/0/o, 0/2.5/p, 0/-2.5/q}
  				\tkzDefShiftPoint[o](-60:2){0}
  				\tkzDefShiftPoint[o](0:2){m'}
  				\tkzDefPointBy[projection=onto m--n](0) \tkzGetPoint{h}
  				\draw(o)circle(2cm);
  				\tkzDrawSegments[dashed](0,h)
  				\tkzDrawSegments(p,q)
  				\tkzDrawSegments[->](m,n)
  				\tkzDrawSegments[blue,->](o,0 o,m')
  				\tkzDrawPoints[size=3](o,h,0,m')
  				\tkzMarkAngles[size=0.7cm,arrows=->](0,o,m')
  				\node at (0)[below right]{$M_0$};
  				\node at (m')[below right]{$M'$};
  				\node at (2,0)[above right]{$10$};
  				\node at (h)[above]{$5$};
  		\end{tikzpicture}}
  		\begin{itemize}
  			\item Tốc độ trung bình của vật là $v_{tb} = \dfrac{45}{\dfrac{7}{3}} = 19,29$ m/s.
  		\end{itemize}
  	}
  	%\haidong{6}
  \end{ex} \haidong{20}
  
  %\loai{Tốc độ trung bình giữa hai thời điểm có vận tốc và gia tốc cho trước}
  \begin{ex}%[3P1KD-4][Loai1] 
  	Một vật nhỏ dao động điều hòa trên trục Ox với biên độ 5 cm và tần số bằng 1 Hz. Trong khoảng thời gian vật có vận tốc nhỏ hơn $5\sqrt{3}\pi$ cm/s và gia tốc lớn hơn $10\pi^2\ cm/s^2$, tốc độ trung bình của vật nhỏ bằng
  	\choice
  	{12, 5 cm/s}
  	{10 cm/s}
  	{\True 15 cm/s}
  	{20 cm/s}
  	\loigiai{
  		\immini{\begin{itemize}
  				\item
  				Ta có:
  				$\omega = 2\pi\ rad/s,\ v_{\max} = 10\pi\ cm/s,\ a_{\max} = 20\pi ^2\ cm/s^2$.
  				\item Sử dụng đường tròn hỗn hợp ta dễ dàng thấy, khoảng thời gian vật có vận tốc nhỏ hơn $5\sqrt{3}\pi$ cm/s và gia tốc lớn hơn $10\pi ^2\ cm/s^2$ ứng với khi vật chuyển động từ vị trí $-\dfrac{A}{2}$ theo chiều âm đến vị trí $-\dfrac{A}{2}$ theo chiều dương.
  				\item Quãng đường vật đi được trong khoảng thời gian đó là $S = A = 5$ cm, hết thời gian $\dfrac{T}{3} = \dfrac{1}{3}$ s.
  				\item Tốc độ trung bình $v_{tb}=\dfrac{5}{\dfrac{1}{3}}= 15$ cm/s.
  		\end{itemize}}{\begin{tikzpicture}[scale=1,thick,>=stealth']
  				\tkzInit[ymin=-3,ymax=3,xmin=-4,xmax=4]\tkzClip
  				\tkzDefPoints{-2.5/0/m, 2.5/0/n, 0/0/o, 0/2.5/p, 0/-2.5/q}
  				\tkzDefShiftPoint[o](120:2){n1}
  				\tkzDefShiftPoint[o](-120:2){m1}
  				\tkzDefShiftPoint[o](-60:2){m2}
  				\tkzDefPointBy[projection=onto p--q](m1) \tkzGetPoint{h}
  				\tkzDefPointBy[projection=onto m--n](n1) \tkzGetPoint{h'}
  				\draw(o)circle(2cm);
  				\tkzDrawSegments[dashed](m1,m2 n1,m1)
  				\tkzDrawSegments[->](m,n p,q n,m)
  				\tkzDrawSegments[blue,->](o,m1 o,m2 o,n1)
  				\tkzDrawPoints[size=3](o,m1,m2,h,h')
  				\tkzMarkAngles[color=red,size=2cm,arrows=->,line width=1.5pt](n1,o,m1)
  				\node at (-2,0)[above left]{$20\pi^2$};
  				\node at (-2,0)[below left]{$a$};
  				\node at (0,-2)[below right]{$10\pi$};
  				\node at (0,-2)[below left]{$v$};
  				\node at (h)[above right]{$5\sqrt{3}\pi$};
  				\node at (h')[above]{$10\pi^2$};
  		\end{tikzpicture}}
  		
  	}
  	%\haidong{6}
  \end{ex} \haidong{20}
 
  
  \setcounter{ex}{0}
  \PhanII
\begin{ex}
  	Vật dao động điều hòa có đồ thị li độ phụ thuộc thời gian như hình vẽ.
  	\immini{\choiceTF[t]
  	{\True Quãng đường vật đi được sau $0,6\ s$ là $12\ cm$}
  	{\True Tại thời điểm $t=0,2\ s$, vật đi qua vị trí cân bằng theo chiều dương}
  	{Tại thời điểm ban đầu, vật ở vị trí biên âm}
  	{Biên độ dao động của vật là $4\ cm$}}{
\begin{hinhanh}{7}
\begin{tikzpicture}[very thick,>=stealth']
\pgfplotsset{width=7cm,height=5cm,compat=1.4}
\begin{axis}[
axis lines=middle, xmin=0, xmax=2*pi, xstep=1,
ymin=-4.5, ymax=5, ystep=0.1,
grid=major,%Vẽ chế độ lưới theo trục tọa độ Oxy,
x grid style={dashed, black},
axis line style={very thick,Mapcolor}, % <--- dòng thêm vào
xlabel=$t~(s)$,
xlabel style={above},
xtick ={1,2,3,4,5,6,7,8},
xticklabels={},
y grid style={dashed, black},
extra y ticks={0},
extra y tick label={0},
ylabel=$x~(cm)$,
ylabel style={above},
ytick={-4,-2,0,2, 4},
yticklabels={$-2$,,,,2},]
\addplot[line width=1.2pt, domain=0:1.7*pi, samples=400, Mapcolor]{4*cos(deg(0.5*pi*x+pi/2))};
\addplot[Mapcolor,mark=*,mark options={fill=white,mark size=1.5pt}, nodes near coords={0,1}]
coordinates {(1, 0)};
\addplot[Mapcolor,mark=*,mark options={fill=white,mark size=1.5pt}, nodes near coords={0,3}]
coordinates {(3, 0)};
\addplot[Mapcolor,mark=*,mark options={fill=white,mark size=1.5pt}, nodes near coords={0,5}]
coordinates {(5, 0)};
\end{axis}
\end{tikzpicture}\end{hinhanh}}
  	\loigiai{
  		\begin{itemchoice}
  			\itemch 
  			\itemch 
  			\itemch 
  			\itemch 
  		\end{itemchoice}
  	}
  \end{ex}
  	\begin{ex}
  	Vật nhỏ dao động điều hòa theo phương trình $x=10\cos\left(2\pi t-\dfrac{2\pi}{3}\right)\ (cm)$. Với $t$ tính bằng giây.
  	\choiceTF[t]
  	{\True Tần số góc của dao động là $2\pi\ rad/s$}
  	{\True Pha dao động của vật lúc $t=2\ s$ là $\dfrac{10\pi}{3}\ rad$}
  	{Gia tốc của vật khi ở biên âm là $-40\pi^{2}\ m/s^{2}$}
  	{\True Trong thời gian $2,5\ s$ vật đi được quãng đường là $100\ cm$}
  	\loigiai{
  		\begin{itemchoice}
  			\itemch
  			\itemch
  			\itemch
  			\itemch
  		\end{itemchoice}
  	}
  \end{ex}
  \begin{ex}
  	Một chất điểm dao động điều hòa theo trục $Ox$, với $O$ trùng với vị trí cân bằng của chất điểm. Đường biểu diễn sự phụ thuộc li độ chất điểm theo thời gian $t$ cho ở hình vẽ.
  	\immini{\choiceTF[t]
  	{\True Khoảng thời gian ngắn nhất giữa 2 lần chất điểm đổi chiều chuyển động là $0,1\ s$}
  	{Biên độ dao động của vật là $12\ cm$}
  	{Tại thời điểm $t=0,1\ s$ li độ của chất điểm là $3\ cm$ và đang chuyển động theo chiều dương}
  	{\True Quãng đường vật đi được sau khoảng thời gian $\Delta t=\dfrac{2}{15}\ s$, kể từ thời điểm ban đầu là $18\ cm$}}{
\begin{hinhanh}{7}
\begin{tikzpicture}[very thick,>=stealth']
\pgfplotsset{width=7cm,height=5cm,compat=1.4}
\begin{axis}[
axis lines=middle, xmin=0, xmax=2.45*pi, xstep=1,
ymin=-4.5, ymax=5, ystep=0.1,
grid=major,%Vẽ chế độ lưới theo trục tọa độ Oxy,
x grid style={dashed, black},
axis line style={very thick,Mapcolor}, % <--- dòng thêm vào
xlabel=$t~(s)$,
xlabel style={above},
xtick ={2/3,4/3,6/3,8/3,10/3,12/3,14/3,16/3,18/3,20/3},
xticklabels={},
y grid style={dashed, black},
extra y ticks={0},
extra y tick label={0},
ylabel=$x~(cm)$,
ylabel style={above},
ytick={-4,-2,0,2, 4},
yticklabels={$-6$,$-3$,,3,6},]
\addplot[line width=1.2pt, domain=0:2.4*pi, samples=400, Mapcolor]{4*cos(deg(0.5*pi*x-2*pi/3))};
\addplot[Mapcolor,mark=*,mark options={fill=white,mark size=1.5pt}, nodes near coords={}]
coordinates {(20/3, -2)};
\addplot[Mapcolor,mark=*,mark options={fill=white,mark size=0.5pt}, nodes near coords={0,2}]
coordinates {(20/3, -4)};
\end{axis}
\end{tikzpicture}
\end{hinhanh}}
  	\loigiai{
  		
  	}
  \end{ex}
  	\begin{ex}
  	Cho một vật dao động điều hoà với phương trình $x=9\cos\left(2\pi t-\dfrac{\pi}{6}\right)\ (cm)$, $t$ tính bằng giây.
  	\choiceTF[t]
  	{\True Quãng đường vật đi được trong một chu kì là $36\ cm$}
  	{\True Phương trình vận tốc của vật là $v=18\pi\cos\left(2\pi t+\dfrac{\pi}{3}\right)\ (cm/s)$}
  	{Gia tốc của vật tại thời điểm $t=2\ s$ là $-18\sqrt{3}\pi\ cm/s^2$}
  	{Từ thời điểm ban đầu đến thời điểm $t=43,5\ s$, vật đi được quãng đường là $391,5\ cm$}
  	\loigiai{
  		\begin{itemchoice}
  			\itemch
  			\itemch
  			\itemch
  			\itemch
  		\end{itemchoice}
  	}
  \end{ex}
  
  \setcounter{ex}{0}
  \PhanIII
    \begin{ex}
    	Một vật dao động điều hòa theo phương trình $x=10 \cos (6,28 t+3,14) cm$. Quãng đường vật đi được sau 3 dao động là bao nhiêu mét?
    	\shortans[oly]{1,2}
    	\loigiai{
    		Phương trình dao động là $x=10 \cos (6,28 t+\pi) cm$.
    		Biên độ dao động của vật là $A = 10\ cm$.
    		Trong một dao động toàn phần, vật đi được quãng đường là $4A = 4 . 10 cm = 40\ cm$.
    		Sau 3 dao động, quãng đường vật đi được là $3 . 4A = 3 . 40 cm = 120\ cm$.
    		Đổi sang mét: $120 cm = 1,2\ m$.
    	}
    \end{ex}
    \begin{ex}
    	Vật dao động điều hòa có phương trình vận tốc là $v=8 \pi \cos \left(2 \pi t-\dfrac{2 \pi}{3}\right)\ cm/s$. Thời gian để vật đi được quãng đường $40\ cm$ kể từ lúc bắt đầu dao động là bao nhiêu giây?
    	\shortans[oly]{2,5}
    	\loigiai{
    	}
    \end{ex}
    \begin{ex}
  	Một chất điểm dao động điều hòa với phương trình li độ $x=10 \cos \left(4 \pi t+\dfrac{\pi}{6}\right)$
  	(t tính bằng giây). Quãng đường chất điểm đi được bằng bao nhiêu $m$ kể từ thời điểm $t=0$ đến thời điểm $t=\dfrac{11}{6} \ s$ (làm tròn kết quả đến chữ số hàng phần trăm)?\\
  	\shortans[oly]{1,49}
  	\loigiai{
  		$1,49 \ m$
  	}
  	%\haidong{6}
  \end{ex} \haidong{20}
  \begin{ex}
  	Một chất điểm dao động điều hòa với phương trình li độ $x=5 \cos \left(4 \pi t-\dfrac{\pi}{3}\right)$ (t tính bằng giây). Tìm thời gian để chất điểm đi được quãng đường $45 \ cm$ kể từ thời điểm $t=0$ (làm tròn kết quả đến chữ số hàng phần trăm)?\\
  	\shortans[oly]{1,17}
  	\loigiai{
  		$1,17 \ s$
  	}
  	%\haidong{6}
  \end{ex} \haidong{20}
  \begin{ex}
  	Một chất điểm dao động điều hòa với phương trình li độ $x=10 \cos \left(\pi t+\dfrac{\pi}{3}\right)$
  	(cm) ($t$ tính bằng giây). Tìm thời gian để chất điểm đi được quãng đường $30 \ cm$ kể từ thời điểm $t=0$ (làm tròn kết quả đến chữ số hàng phần trăm)?\\
  	\shortans[oly]{1,33}
  	\loigiai{
  		$1,33 \ s$
  	}
  	%\haidong{6}
  \end{ex} \haidong{20}
  \begin{ex}
  	Một chất điểm dao động điều hòa với phương trình li độ $x=5 \cos (4 \pi t+\pi)\ cm$ (t tính bằng giây). Tìm thời gian để chất điểm đi được quãng đường $255 \ cm$ kể từ thời điểm $t=\dfrac{1}{6} \ s$ (làm tròn kết quả đến chữ số hàng phần trăm)?\\
  	\shortans[oly]{6,42}
  	\loigiai{
  		$6,42 \ s$
  	}
  	%\haidong{6}
  \end{ex} \haidong{20}
  
  
  
  
  